\documentclass{abntex2}
\usepackage[utf8]{inputenc}
\usepackage[T1]{fontenc}
\usepackage{geometry}
\usepackage{graphicx}
\usepackage{float}
\usepackage{hyperref}
\usepackage[english,indonesian]{babel}

\title{Optimasi Pemrosesan Template LaTeX Menggunakan Pendekatan Berbasis AI}
\author{Academic Support Agent}
\date{\today}

\begin{document}
\maketitle

\begin{abstract}

\end{abstract}

\chapter{PENDAHULUAN}


\chapter{TINJAUAN PUSTAKA}


% Original content (commented out):
% \documentclass[12pt,a4paper,oneside]{abntex2}

\input{configuracoes}

% Informações de dados para CAPA e FOLHA DE ROSTO
\title{Título da Dissertação}
\author{Nome do Autor}
\orientador{Prof. Dr. Nome do Orientador}
\instituicao{UNIVERSIDADE FEDERAL DO ESPÍRITO SANTO \par PROGRAMA DE PÓS-GRADUAÇÃO EM MATEMÁTICA}
\date{2023}
\local{Vitória}
\tipotrabalho{Dissertação de Mestrado}
\preambulo{Dissertação de mestrado apresentada
ao PPGMAT como parte dos requisitos exigidos para a obtenção do título de Mestre em Matemática}

\makeindex

\begin{document}

% Seleciona o idioma do documento 
\selectlanguage{brazil}

% Capa
\imprimircapa

% Folha de rosto
\imprimirfolhaderosto*

% Para gerar a ficha bibliográfica acesse:
% https://biblioteca.ufes.br/ficha-catalografica-cip

% Inserir a ficha bibliografica gerada pela biblioteca
\begin{fichacatalografica}
    \includepdf{fig_ficha_catalografica.pdf}
\end{fichacatalografica}

% Folha de aprovação
\begin{folhadeaprovacao}
\includepdf{folhadeaprovacao_final.pdf}
\end{folhadeaprovacao}

% Folha de aprovação editável
\begin{folhadeaprovacao}

  \begin{center}
    % Remover após edição da folha de aprovação
    \textbf{Editável}
    
    {\ABNTEXchapterfont\large\imprimirautor}

    \vspace*{\fill}\vspace*{\fill}
    \begin{center}
      \ABNTEXchapterfont\bfseries\Large\imprimirtitulo
    \end{center}
    \vspace*{\fill}
    
    \hspace{.45\textwidth}
    \begin{minipage}{.5\textwidth}
        \imprimirpreambulo
    \end{minipage}%
    \vspace*{\fill}
   \end{center}
        
   Trabalho aprovado. \imprimirlocal, \today:

   \assinatura{\textbf{\imprimirorientador} \\ Universidade Federal do Espírito Santo \\ Orientador} 
   \assinatura{\textbf{Prof. Dr. Convidado 1} \\ Universidade Federal do Espírito Santo \\ Membro Interno}
   \assinatura{\textbf{Prof. Dr. Convidado 2} \\ Instituição \\ Membro Externo}
   
   \begin{center}
    \vspace*{0.5cm}
    {\large\imprimirlocal}
    \par
    {\large\imprimirdata}
    \vspace*{1cm}
  \end{center}
  
\end{folhadeaprovacao}


% Dedicatória
\begin{dedicatoria}
   \vspace*{\fill}
   \centering
   \noindent
   \textit{Texto referente a dedicatória.} \vspace*{\fill}
\end{dedicatoria}


% Agradecimentos
\begin{agradecimentos}
Texto referente aos agradecimentos.
\end{agradecimentos}


% Epígrafe
\begin{epigrafe}
    \vspace*{\fill}
	\begin{flushright}
		\textit{``Texto referente a epigrafe." \\
		(Indicar fonte do texto)}
	\end{flushright}
\end{epigrafe}


% RESUMOS

% Resumo em português
\setlength{\absparsep}{18pt} % Espaçamento dos parágrafos
\begin{resumo}
Texto referente ao resumo em português.

 \textbf{Palavras-chave}: Palavras-chave em português
\end{resumo}

% resumo em inglês
\begin{resumo}[Abstract]
 \begin{otherlanguage*}{english}
   This is the english abstract.

   \vspace{\onelineskip}
 
   \noindent 
   \textbf{Keywords}: Keywords in english
 \end{otherlanguage*}
\end{resumo}


% Lista de ilustrações
\pdfbookmark[0]{\listfigurename}{lof}
\listoffigures*
\cleardoublepage

% Lista de quadros
%\pdfbookmark[0]{\listofquadrosname}{loq}
%\listofquadros*
%\cleardoublepage

% Lista de tabelas
\pdfbookmark[0]{\listtablename}{lot}
\listoftables*
\cleardoublepage

% Lista de abreviaturas e siglas
%\begin{siglas}
%  \item[ABNT] Associação Brasileira de Normas Técnicas
%  \item[abnTeX] ABsurdas Normas para TeX
%\end{siglas}

% Lista de símbolos
%\begin{simbolos}
%  \item[$ \Gamma $] Letra grega Gama
%  \item[$ \Lambda $] Lambda
%  \item[$ \zeta $] Letra grega minúscula zeta
%  \item[$ \in $] Pertence
%\end{simbolos}

% Sumário
\pdfbookmark[0]{\contentsname}{toc}
\tableofcontents*
\cleardoublepage


% ---------------------------------------------------------
% ELEMENTOS TEXTUAIS
\textual

% ---------------------------------------------------------
\chapter{Introdução}

Texto referente à introdução.

% ---------------------------------------------------------
\chapter{Conteúdos específicos}
\label{cap:conteudos_especificos}

\section{Definições, Teoremas, Corolários e Exemplos}

A Definição \ref{def:modelo_definicao} mostra um modelo de definição.

\begin{definicao}
Aqui temos uma nova definição.
\label{def:modelo_definicao}
\end{definicao}

O Teorema \ref{teo:modelo_teorema} mostra um modelo de teorema.

\begin{teorema}
Aqui temos um novo teorema.
\label{teo:modelo_teorema}
\end{teorema}

O Exemplo \ref{ex:modelo_exemplo} mostra um modelo de exemplo.

\begin{exemplo}
Aqui temos um novo exemplo.
\label{ex:modelo_exemplo}
\end{exemplo}

O Corolário \ref{cor:modelo_corolario} mostra um modelo de corolário.

\begin{corolario}
Aqui temos um novo corolário.
\label{cor:modelo_corolario}
\end{corolario}


\section{Citações das Referências}

Livro de Análise Real~\cite{lima2004analise}


\section{Listas}

Exemplo de lista com símbolos:

\begin{itemize}
    \item Primeiro elemento
    \item Segundo elemento
\end{itemize}

Exemplo de lista com itens enumerados:

\begin{enumerate}
    \item Primeiro elemento
    \item Segundo elemento
\end{enumerate}


\section{Tabelas}

A Tabela~\ref{tab:modelo_tabela} mostra um modelo de tabela.

\begin{table}[htb]
    \caption{Exemplo de tabela}
    \label{tab:modelo_tabela}
    \centering
    \begin{tabular}{c|c|c|c}
	    \hline
	    \textbf{Pessoa} & \textbf{Idade} & \textbf{Peso} & \textbf{Altura} \\ \hline
	    Marcos & 26    & 68   & 178    \\ \hline
	    Ivone  & 22    & 57   & 162    \\ \hline
	    Sueli  & 40    & 65   & 153    \\ \hline
    \end{tabular}
    \fonte{Produção do próprio autor (2023).}
\end{table}


\section{Imagens}

A Figura~\ref{fig:logo_UFES} mostra a logomarca promocional da Universidade Federal do Espírito Santo (UFES).

\begin{figure}[htb]
    \caption{Logomarca Promocional UFES}
    \label{fig:logo_UFES}
    \centering
    \includegraphics[width=0.5\textwidth]{imagens/logomarca_UFES}
    \fonte{Produção do próprio autor (2023).}
\end{figure}


% ---------------------------------------------------------
\phantompart

% Conclusão
\chapter{Conclusão}


Texto referente à conclusão.

% ---------------------------------------------------------
% ELEMENTOS PÓS-TEXTUAIS
\postextual

% ---------------------------------------------------------% Referências bibliográficas

\bibliography{referencias}


% Inicia os apêndices
\begin{apendicesenv}

%\partapendices

\chapter{Título do Apêndice}

Texto referente ao apêndice.

%\chapter{Texto referente ao apêndice.}
%
%Texto referente ao apêndice.

\end{apendicesenv}

\end{document}


\end{document}